% TOP_PROBLEM: {"problemId": "problem.nagata-conjecture-on-plane-curve-linear-systems", "canonicalTitle": "Nagata Conjecture on Plane Curve Linear Systems", "releaseRank": 365, "primaryDomain": "geometry_topology", "primaryDomainLabel": "Geometry & topology", "displayStatus": "open", "releaseStatus": "open"}
% !TeX program = lualatex
% Definition and sources only; this file is not an LLM solution attempt.
\documentclass[11pt]{article}
\usepackage[margin=25mm]{geometry}
\usepackage{fontspec}
\setmainfont{Latin Modern Roman}
\usepackage{amsmath,amssymb,mathtools}
\usepackage{unicode-math}
\setmathfont{Latin Modern Math}
\usepackage{xurl,hyperref}
\hypersetup{colorlinks=true,urlcolor=blue}
\setcounter{secnumdepth}{0}
\setlength{\parindent}{0pt}
\setlength{\parskip}{5pt}
\setlength{\emergencystretch}{3em}
\begin{document}
\section*{\texorpdfstring{Nagata Conjecture on Plane Curve Linear Systems}{Nagata Conjecture on Plane Curve Linear Systems}}\label{365}
\subsection{Definitions and mathematical statement}
Choose $r\ge10$ very general points $p_1,\ldots,p_r\in{\mathbb{P}}^2_{{\mathbb{C}}}$, meaning outside a countable union of proper closed subsets of the configuration space. For a plane curve $C$ of positive degree $d$, its multiplicity at $p_i$ is the vanishing order of a local defining equation. The selected Nagata inequality asserts that, whenever $\operatorname{mult}_{p_i}C\ge m_i$ for integers $m_i\ge0$,
\[
 d>\frac1{\sqrt r}\sum_{i=1}^r m_i.
\]
The configuration is very general, not every arrangement of distinct points. Unequal multiplicities and reducible curves are included in the statement. The restriction to actual positive-degree curves excludes the empty degree-zero divisor, for which a strict $0>0$ inequality would be meaningless.
\par\smallskip\textit{Formulation and concept references:} \href{https://link.springer.com/article/10.1007/s40687-025-00500-2}{[S1]}.

\subsection{Short English statement}
For sufficiently many very general points in the projective plane, must a curve's degree strictly exceed the stated average multiplicity bound?
\subsection{Sources}
\begin{itemize}
\item[S1]Carlos Galindo et al., On the valuative Nagata conjecture, Research in the Mathematical Sciences 12, 18 (2025).
\url{https://link.springer.com/article/10.1007/s40687-025-00500-2}.
\textit{Formulation source.}
\end{itemize}
\end{document}
